% CP-VI-0037
% target_chapter: VI/27 — Integral Schemes and Function Fields
% source_unit_id: IMP-DL-D27460A55B-SPR003063-B0D7BDB0
% source: Downloads/theory-of-algebraic-geometry-5.tex L3063-L3071
% solution_provenance: CANONICAL_AUTHORED_SOLUTION

\textbf{Problem.}
Let \(X\) be an integral scheme. Show that there is a unique point \(\eta\) such that the closure of \(\{\eta\}\) is \(X\); this is called the \emph{generic point} of \(X\). Show that the stalk of \(\mathcal{O}_X\) at \(\eta\) is a field, called the \emph{function field} of \(X\), denoted by \(K(X)\). Show that if
\[
U=\operatorname{Spec} A
\]
is any open affine subset of \(X\), then \(K(X)\) is the field of fractions of \(A\).

\bigskip

\subsection*{Canonical solution}
Because $X$ is irreducible, every nonempty affine open $U=\operatorname{Spec} A$ is irreducible; since $X$ is also reduced, $A$ is a domain. The generic point of $U$ is therefore the prime $(0)$. If $U$ and $V$ are two nonempty affine opens, then $U\cap V\neq\varnothing$, and the generic points of $U$ and $V$ agree on the intersection. They consequently define a unique point $\eta\in X$ whose closure is all of $X$.

For $U=\operatorname{Spec} A$ containing $\eta$, the stalk is
\[
\mathcal O_{X,\eta}\cong A_{(0)}=\operatorname{Frac}(A),
\]
a field. Every nonempty affine open contains $\eta$, so the same stalk computes the fraction field of the coordinate ring of any such affine open. Thus $K(X)=\mathcal O_{X,\eta}\cong\operatorname{Frac}(A)$.
